\documentclass{article}

\usepackage{amsmath}
\usepackage{amssymb}
\usepackage{amsthm}
\usepackage[top=1in, bottom=1in, left=1.25in, right=1.25in]{geometry}
\usepackage{graphicx}
\usepackage{minted}
\usepackage{mathrsfs}
\usepackage[shortlabels]{enumitem}

\theoremstyle{plain}
\newtheorem{thm}{Theorem}
\newtheorem{lem}[thm]{Lemma}
\newtheorem{prop}[thm]{Proposition}
\newtheorem{cor}[thm]{Corollary}
\newtheorem{problem}{Problem}

% shortcuts for blackboard bold number sets (reals, integers, etc.)
\newcommand{\QQ}{\ensuremath{\mathbb Q}}
\newcommand{\RR}{\ensuremath{\mathbb R}}
\newcommand{\NN}{\ensuremath{\mathbb N}}
\newcommand{\ZZ}{\ensuremath{\mathbb Z}}

% Make header with name and date etc.
\usepackage{fancyhdr}
\lhead{Kaelin Ellis\\Math 320: Combinatorics}
\rhead{\today\\Assignment 1}
\pagestyle{fancy}

\usepackage[utf8]{inputenc}
\setlength{\parindent}{0pt} % Don't indent new paragraphs
\setlength{\headheight}{24pt}

\setcounter{problem}{0}

\begin{document}

\section{Chapter 1.1}

\begin{problem} % Problem 1
How many different tickets are possible in each of the following lotteries? And which lottery offers the best chance of winning?
\begin{enumerate}[(a)]
    \item You pick six numbers from 1-16, a number can be picked more than once, and order doesn't matter.
    \item You pick five different numbers from 1-25 and order doesn't matter.
    \item You pick four different numbers from 1-18 and the order in which you specify them matters.
\end{enumerate}
\end{problem}

\textbf{Answer:}
\begin{enumerate}[(a)]
    \item A multiset combination where $k = 6$ and $n=16$ which is:
    \begin{equation*}
        \dbinom{6+16-1}{6} = \dbinom{21}{6} = \dfrac{21!}{6!15!} = 54,264
    \end{equation*}
    \item A subset combination where $k = 5$ and $n=25$ which is:
    \begin{equation*}
        \dbinom{25}{5} = \dfrac{25!}{5!(25-5)!} = \dfrac{25!}{5!20!} = 53,130
    \end{equation*}
    \item A permutation with no repetition where $k = 4$ and $n=18$.
    \begin{equation*}
        (18)_{4} = \dfrac{18!}{(18-4)!} = \dfrac{18!}{14!} = 73,440
    \end{equation*}
\end{enumerate}
The second lottery appears to have the best odds of winning.

\begin{problem} % Problem 2

You flip a coin 20 times and record the ordered sequence of heads and tails.

\begin{enumerate}[(a)]
\item How many sequences are there in which you get heads on (at least) flip \#1, \#4, \#7, and \#13?

\item How many sequences have the same number of heads and tails?
\end{enumerate}
\end{problem}

\textbf{Answer:}
\begin{enumerate}[(a)]
    \item Each choice can be a $\{0, 1\}$ for heads or tails. As a result $n=2$. The positions 1, 4, 7, and 13 have already been decided. This leaves 16 positions available for a head or tails selection. The answer is $2^{16}$.
    \item How many combinations of 20 positions have exactly 10 heads? The answer is $\dbinom{20}{10}$. How many times can you write a 20 character long string with exactly 10 digits being locked?
\end{enumerate}


\setcounter{problem}{8}
\begin{problem} % Problem 9
How many subsets of $[20] \dots$
\begin{enumerate}[(a)]
\item Have smallest element 4 and largest element 15?
\item Contain no even numbers?
\item Have size 10 and don't contain any number larger than 17?
\end{enumerate}
\end{problem}

\textbf{Answer:}
\begin{enumerate}[(a)]
    \item How many subsets can be written with the numbers between 4 and 15, do not include the empty set. If $B = \{x\in \ZZ: 4 \le x \le 15 \} = \{4, 5, 6, \dots 15\}$. This is asking for the size of the power set of B, but ignoring the empty set. $|\mathscr{P}(B)| = 2^{15-3} = 4096$. The answer is $4096-1 = 4095$. This problem can be thought as binary sequence of 12 digits. Each digit corresponds to whether the number 4 - 15 is included. If the digit is a one, it is included. If the digit is zero, it is not included.
    \item Again, using the analogy that this problem can be thought as a binary sequence with a 0 implying the digit is not included and a 1 implying the digit is included. As such, all the even numbers are automatically swapped to a 0, so only 1 choice. The odd numbers can be a 0 or 1, so two choices. There are 10 odd numbers between 1 and 20. So the answer is $2^{10}$.
    \item This is asking for a subset that is length 10, or $k = 10$. The number of choices has been restricted to 17, so $n=17$. It is a subset, so no repetition is allowed. As such the answer is $\dbinom{17}{10}$.
\end{enumerate}

\setcounter{problem}{13}
\begin{problem} % Problem 14
How many arthmetic problems of the following form are possible? You must use each of the digits 1 through 9, they must appear in numerical order from left to right, and you can use any combination of $+$ and $\times$ symbols you like, as long as the resulting expression makes mathematical sense. For example $1234 + 5 \times 6 \times 78 + 9$ and $123456 + 789$ and $123456789$ are three possibilities, but $1 \times \times 234567 + 89$ is not.
\end{problem}

\textbf{Answer:}\smallskip

In between the sets of two numbers are eight spaces. Each one of these spaces can be inserted with no symbol, a $+$ symbol or a $\times$ symbol. This means there are 3 choices or $n=3$ for a list of length 8, or $k=8$. The answer is $3^8$.


\setcounter{problem}{16}
\begin{problem} % Problem 17
Find the number of 3-lists of the form $(x_1, x_2, x_3)$, where each $x_i$ is a nonnegative integer and $x_1 + x_2 + x_3 = 10$.
\end{problem}

\textbf{Answer:}\smallskip

Define $X = \{x_1, x_2, x_3\}$, where $X$ is a multiset. If the values in $X$ are nonnegative and all 3 values must sum to 10, then $X \subseteq A$. Where
\begin{equation*}
    A = \{0, 1, 2, 3, 4, 5, 6, 7, 8, 9, 10\}
\end{equation*}
Because this is a list, the $x_1$ sets the appropriate range of the $x_2$ with $x_2 \le 10 - x_1$. Then $x_1$ and $x_2$ set $x_3$ with $x_3 = 10 - x_1 - x_2$. As such, only the first two values have multiple options.
\begin{enumerate}
    \item If $x_1 = 0$ then $x_2 \le 10$, which has 11 choices.
    \item If $x_1 = 1$ then $x_2 \le 9$, which has 10 choices.
    \item If $x_1 = 2$ then $x_2 \le 8$, which has 9 choices.
    \item If $x_1 = 3$ then $x_2 \le 7$, which has 8 choices.
    \item $\dots$
    \item If $x_1 = 10$ then $x_2 = 0$, which is 1 choice.
\end{enumerate}
As the value for $x_1$ increase, the options for $x_2$ decreases. The answer is represented by $11!$.
\end{document}

